\chapter*{Course Map: From Search to Intelligent Systems}
\addcontentsline{toc}{chapter}{Course Map: From Search to Intelligent Systems}
\markboth{Course Map}{}

The eight laboratory handouts form a compact route through several classical Artificial Intelligence ideas. The course begins with search, then moves to approximate reasoning and control, language processing, and evolutionary optimization. Use this map to see how each lab contributes to the whole.

\begin{mlfigure}
\begin{tikzpicture}[
  stage/.style={draw=MLBlue,fill=MLPaleBlue,rounded corners=2mm,text width=12.7cm,minimum height=1.45cm,align=left,inner xsep=4mm,inner ysep=2.5mm,font=\scriptsize\color{MLNavy}},
  stage2/.style={draw=MLTeal,fill=MLPaleTeal,rounded corners=2mm,text width=12.7cm,minimum height=1.45cm,align=left,inner xsep=4mm,inner ysep=2.5mm,font=\scriptsize\color{MLNavy}},
  stage3/.style={draw=MLOrange,fill=MLPaleOrange,rounded corners=2mm,text width=12.7cm,minimum height=1.45cm,align=left,inner xsep=4mm,inner ysep=2.5mm,font=\scriptsize\color{MLNavy}},
  stage4/.style={draw=MLGreen,fill=MLPaleTeal,rounded corners=2mm,text width=12.7cm,minimum height=1.45cm,align=left,inner xsep=4mm,inner ysep=2.5mm,font=\scriptsize\color{MLNavy}},
  arr/.style={-{Stealth[length=2mm]},thick,draw=MLMuted}
]
\node[stage] (s1) at (0,4.8) {{\normalsize\bfseries\color{MLBlue}Stage 1: Search and problem solving}\hfill{\footnotesize\bfseries Labs 1--3}\par
\smallskip\textbf{Learn:} BFS, DFS, Uniform-Cost Search, Greedy Best-First Search, A*, grids, graphs, trees, and word ladders.\par
\textbf{Outcome:} choose a frontier rule that matches the information available.};
\node[stage2] (s2) at (0,2.5) {{\normalsize\bfseries\color{MLTeal}Stage 2: Approximate reasoning and control}\hfill{\footnotesize\bfseries Labs 4--5}\par
\smallskip\textbf{Learn:} fuzzy membership, FIS, ANFIS, robot motion, PID line following, path planning, collision avoidance, and forward kinematics.\par
\textbf{Outcome:} connect reasoning rules to simulated intelligent behavior.};
\node[stage3] (s3) at (0,0.2) {{\normalsize\bfseries\color{MLOrange}Stage 3: Language intelligence}\hfill{\footnotesize\bfseries Labs 6--7}\par
\smallskip\textbf{Learn:} embeddings, similarity, OCR, simple captioning, NER simulation, summarization, sentiment, language detection, and keyword extraction.\par
\textbf{Outcome:} represent and process human language with interpretable pipelines.};
\node[stage4] (s4) at (0,-2.1) {{\normalsize\bfseries\color{MLGreen}Stage 4: Evolutionary optimization}\hfill{\footnotesize\bfseries Lab 8}\par
\smallskip\textbf{Learn:} population, chromosome, fitness, selection, crossover, mutation, and iterative optimization.\par
\textbf{Outcome:} formulate a search problem as evolution over candidate solutions.};
\draw[arr] (s1) -- (s2);
\draw[arr] (s2) -- (s3);
\draw[arr] (s3) -- (s4);
\end{tikzpicture}
\end{mlfigure}

\begin{keyidea}
The course does not use one single definition of ``intelligence.'' Instead, it studies several computational abilities: finding paths, using heuristic knowledge, reasoning with degrees of truth, controlling robots, processing language, and optimizing candidate solutions.
\end{keyidea}
